% LaTeX Template for AI Problem Analysis Output (v2.0)
% Strategic Analysis of AMC12B 2023 Problem 16
% IMPORTANT: Use LaTeX commands for all mathematical symbols, NOT unicode characters

\documentclass[12pt]{article}
\usepackage[utf8]{inputenc}
\usepackage{amsmath, amssymb, amsthm}
\usepackage{geometry}
\usepackage{xcolor}
\usepackage[most]{tcolorbox}
\usepackage{enumitem}
\usepackage{fancyhdr}

% Page setup
\geometry{margin=1in}
\setlength{\headheight}{14.5pt}
\pagestyle{fancy}
\fancyhf{}
\rhead{AMC12B 2023 Problem 16 Analysis}
\lhead{\today}

% Color definitions
\definecolor{problemconfig}{RGB}{230, 242, 255}    % Light blue
\definecolor{strategic}{RGB}{255, 230, 242}       % Light pink
\definecolor{tactical}{RGB}{242, 255, 230}        % Light green
\definecolor{techniques}{RGB}{255, 242, 230}      % Light yellow
\definecolor{verification}{RGB}{255, 230, 230}    % Light red
\definecolor{solution}{RGB}{230, 255, 255}        % Light cyan
\definecolor{competition}{RGB}{242, 242, 255}     % Light purple
\definecolor{proofwriting}{RGB}{240, 230, 255}    % Light lavender
\definecolor{gold}{RGB}{218, 165, 32}

% Box styles
\newtcolorbox{problemconfigbox}{
    colback=problemconfig,
    colframe=blue,
    title=Problem Configuration Analysis,
    fonttitle=\bfseries,
    arc=3pt,
    boxrule=1pt,
    breakable
}

\newtcolorbox{strategicbox}{
    colback=strategic,
    colframe=purple,
    title=Strategic Framework Application,
    fonttitle=\bfseries,
    arc=3pt,
    boxrule=1pt,
    breakable
}

\newtcolorbox{tacticalbox}{
    colback=tactical,
    colframe=green,
    title=Tactical Methodology Selection,
    fonttitle=\bfseries,
    arc=3pt,
    boxrule=1pt,
    breakable
}

\newtcolorbox{techniquesbox}{
    colback=techniques,
    colframe=orange,
    title=Conversion Techniques Application,
    fonttitle=\bfseries,
    arc=3pt,
    boxrule=1pt,
    breakable
}

\newtcolorbox{verificationbox}{
    colback=verification,
    colframe=red,
    title=Verification Strategy Design,
    fonttitle=\bfseries,
    arc=3pt,
    boxrule=1pt,
    breakable
}

\newtcolorbox{solutionbox}{
    colback=solution,
    colframe=cyan,
    title=Solution Roadmap,
    fonttitle=\bfseries,
    arc=3pt,
    boxrule=1pt,
    breakable
}

\newtcolorbox{competitionbox}{
    colback=competition,
    colframe=purple,
    title=Competition Strategy Integration,
    fonttitle=\bfseries,
    arc=3pt,
    boxrule=1pt,
    breakable
}

\newtcolorbox{proofbox}{
    colback=proofwriting,
    colframe=violet,
    title=Proof Structure and Justification (COMC Part C),
    fonttitle=\bfseries,
    arc=3pt,
    boxrule=1pt,
    breakable
}

\newtcolorbox{keyinsightbox}{
    colback=yellow!20,
    colframe=gold,
    title=Key Insight,
    fonttitle=\bfseries,
    arc=3pt,
    boxrule=2pt,
    leftrule=5pt,
    breakable
}

\newtcolorbox{warningbox}{
    colback=red!10,
    colframe=red!50,
    title=Warning: Common Pitfall,
    fonttitle=\bfseries,
    arc=3pt,
    boxrule=2pt,
    leftrule=5pt,
    breakable
}

\newtcolorbox{tipbox}{
    colback=green!10,
    colframe=green!50,
    title=Pro Tip,
    fonttitle=\bfseries,
    arc=3pt,
    boxrule=2pt,
    leftrule=5pt,
    breakable
}

\begin{document}

\title{AMC12B 2023 Problem 16\\Strategic Analysis}
\author{Competition Math Framework v2.1}
\date{\today}
\maketitle

\section*{Problem Statement}

In the state of Coinland, coins have values $6, 10,$ and $15$ cents. Suppose $x$ is the value in cents of the most expensive item in Coinland that cannot be purchased using these coins with exact change. What is the sum of the digits of $x$?

$$\textbf{(A) } 8 \qquad \textbf{(B) } 10 \qquad \textbf{(C) } 7 \qquad \textbf{(D) } 11 \qquad \textbf{(E) } 9$$

\begin{problemconfigbox}
\subsection*{A. Problem Classification}
\begin{itemize}[leftmargin=*]
    \item \textbf{Problem Type}: To Find (determine the largest non-representable value, then compute digit sum)
    \item \textbf{Competition Context}: AMC 12B 2023, Problem 16
    \item \textbf{Difficulty Band}: Mid-band (P16) - Upper end of mid-band, approaching late-band difficulty
    \item \textbf{Mathematical Domain}: Number Theory with Combinatorial aspects (Frobenius coin problem)
    \item \textbf{Estimated Time}: 4-5 minutes for mid-band P16 (with proper recognition of Frobenius/McNugget theorem)
\end{itemize}

\subsection*{B. Problem Structure Identification}
\begin{itemize}[leftmargin=*]
    \item \textbf{Given Information}: 
    \begin{itemize}
        \item Three coin denominations: 6, 10, and 15 cents
        \item We need exact change (non-negative integer combinations)
        \item Looking for the \textit{largest} value that \textit{cannot} be represented
    \end{itemize}
    \item \textbf{Constraints}: 
    \begin{itemize}
        \item Non-negative integer coefficients only (can use 0, 1, 2, \ldots of each coin)
        \item Must find the maximum non-representable value
        \item Need to verify that all larger values \textit{are} representable
    \end{itemize}
    \item \textbf{Objective}: 
    \begin{itemize}
        \item Find largest $x$ such that $6a + 10b + 15c = x$ has no solution with $a, b, c \in \mathbb{Z}_{\geq 0}$
        \item Compute the sum of digits of $x$
    \end{itemize}
    \item \textbf{Hidden Assumptions}: 
    \begin{itemize}
        \item Such a maximum value exists (it does, by Frobenius theory)
        \item All sufficiently large values are representable
        \item $\gcd(6, 10, 15) = 1$ is necessary for finite Frobenius number (but here $\gcd = 1$, so we're safe)
    \end{itemize}
    \item \textbf{Answer Choice Leverage}: 
    \begin{itemize}
        \item Answer choices are small single-digit or low two-digit numbers (7-11)
        \item This suggests $x$ is likely in the range 19-99 (since digit sums in that range can be 7-11)
        \item Most likely candidates for $x$: 29 (digit sum 11), 38 (digit sum 11), 47 (digit sum 11), etc.
    \end{itemize}
    \item \textbf{Proof Requirements}: Not a proof problem; computational verification required
\end{itemize}

\subsection*{C. Strategic Orientation}
\begin{itemize}[leftmargin=*]
    \item \textbf{Hypothesis}: This is a Frobenius coin problem (also known as Chicken McNugget theorem)
    \item \textbf{Conclusion}: Need to find the largest non-representable value for the system $6a + 10b + 15c$ with $a, b, c \geq 0$
    \item \textbf{Notation}: 
    \begin{itemize}
        \item Let $a$ = number of 6-cent coins
        \item Let $b$ = number of 10-cent coins
        \item Let $c$ = number of 15-cent coins
        \item Target: $x$ such that $6a + 10b + 15c = x$ has no non-negative integer solutions
    \end{itemize}
    \item \textbf{Intuitive Approach}: 
    \begin{itemize}
        \item Since $\gcd(6, 10) = 2$, we can only make even numbers with 6 and 10
        \item To make odd numbers, we need at least one 15-cent coin
        \item The Frobenius number for two coprime integers $p, q$ is $pq - p - q$
        \item We might be able to reduce this to a two-coin problem plus adjustments
    \end{itemize}
    \item \textbf{Cross-Domain Potential}: 
    \begin{itemize}
        \item Can be viewed through modular arithmetic (parity, divisibility)
        \item Can be approached via systematic enumeration
        \item Can use Chicken McNugget theorem for two coprime values
        \item Visual approach: tabular method (arrange by residue classes)
    \end{itemize}
\end{itemize}
\end{problemconfigbox}

\begin{strategicbox}
\subsection*{A. Psychological Strategies}
\begin{itemize}[leftmargin=*]
    \item \textbf{Mental Toughness}: This problem requires persistence if not immediately recognized as Frobenius problem. If initial approach doesn't work (e.g., trying small cases), pivot to systematic methods.
    \item \textbf{Creativity Cultivation}: Multiple solution paths exist:
    \begin{itemize}
        \item Modular arithmetic approach (considering residues mod 2, mod 3, mod 5, mod 6)
        \item Chicken McNugget theorem applied twice
        \item Visual tabular method
        \item Systematic case analysis by parity
    \end{itemize}
    \item \textbf{Iterative Refinement}: Start with observation that 6 and 10 make only even numbers, so 15 is needed for odd numbers. This constrains the problem significantly.
\end{itemize}

\subsection*{B. Investigation Strategies}
\begin{itemize}[leftmargin=*]
    \item \textbf{Startup Orientation}: 
    \begin{itemize}
        \item Recognize this as a linear Diophantine equation with non-negative constraints
        \item Identify as Frobenius coin problem (Chicken McNugget problem)
        \item Note that $\gcd(6, 10, 15) = 1$, so eventually all sufficiently large integers are representable
    \end{itemize}
    \item \textbf{Get Your Hands Dirty}: 
    \begin{itemize}
        \item Test small values: 1, 2, 3, 4, 5 (none representable initially)
        \item Check which values can be made: 6, 10, 12, 15, 16, 18, 20, 21, \ldots
        \item Identify pattern: odd numbers require at least one 15-cent coin
    \end{itemize}
    \item \textbf{Penultimate Step}: Once we determine that 29 cannot be represented, verify that 30 and all larger values can be represented.
    \item \textbf{Make It Easier}: 
    \begin{itemize}
        \item Simplify by factoring: $6a + 10b + 15c = 2(3a + 5b) + 15c$
        \item Observe parity: even vs. odd numbers require different approaches
        \item Consider the two-coin problem first: $\gcd(6, 10) = 2$, so we can make all even numbers $\geq 2 \cdot (3 \cdot 5 - 3 - 5) = 2 \cdot 7 = 14$ using just 6 and 10
    \end{itemize}
    \item \textbf{Conjecture via Small Cases}: 
    \begin{itemize}
        \item Small odd numbers that cannot be made: 1, 3, 5, 7, 9, 11, 13, 17, 19, 23, 29
        \item Small even numbers that cannot be made: 2, 4, 8
        \item Pattern emerges: 29 appears to be the last gap
    \end{itemize}
    \item \textbf{Visualization}: 
    \begin{itemize}
        \item Create a table modulo 6 to track which residue classes can be achieved
        \item Or create a table showing which values in each column (mod 6) can be reached
    \end{itemize}
\end{itemize}

\subsection*{C. Late-Band Strategic Playbook}
\begin{itemize}[leftmargin=*]
    \item \textbf{Answer-Choice Leverage}: 
    \begin{itemize}
        \item Digit sums 7-11 suggest $x$ is likely in the range 16-99
        \item Working backwards: if digit sum is 11, possible values are 29, 38, 47, 56, 65, 74, 83, 92
        \item Test 29: is it representable? No! (as we'll verify)
    \end{itemize}
    \item \textbf{Make It Easier}: Reduce three-coin problem to two-coin problems
    \item \textbf{Penultimate Step}: Verify that all values $\geq 30$ are representable
    \item \textbf{Wishful Thinking}: If we could show all even numbers $\geq 16$ and all odd numbers $\geq 31$ are representable, then 29 would be the answer
    \item \textbf{Conjecture via Small Cases}: Enumerate non-representable values up to 30, identifying 29 as the largest
\end{itemize}

\subsection*{D. Argument Construction Strategy}
\begin{itemize}[leftmargin=*]
    \item \textbf{Direct Proof}: 
    \begin{enumerate}
        \item Show that all even numbers $\geq 16$ can be represented using only 6 and 10
        \item Show that all odd numbers $\geq 31$ can be represented using one 15 plus 6 and 10
        \item Therefore, all numbers $\geq 30$ are representable
        \item Verify that 29 cannot be represented
        \item Conclude that 29 is the Frobenius number
    \end{enumerate}
    \item \textbf{Stylistic Conventions}: Use Chicken McNugget theorem for two coprime integers
    \item \textbf{Case Analysis}: Split by parity (even vs. odd)
\end{itemize}

\subsection*{E. Cross-Domain Translation}
\begin{itemize}[leftmargin=*]
    \item \textbf{Draw a Picture}: Tabular visualization with rows representing multiples of 6
    \item \textbf{Recast the Problem}: Instead of three coins, think of it as ``even + 15c'' where even numbers come from $6a + 10b$
    \item \textbf{Change Your Point of View}: View through modular arithmetic lens (mod 2, mod 3, mod 5, or mod 6)
\end{itemize}
\end{strategicbox}

\begin{tacticalbox}
\subsection*{A. Core Mathematical Tactics}
\begin{itemize}[leftmargin=*]
    \item \textbf{Parity Arguments}: Central to this problem
    \begin{itemize}
        \item $6a + 10b$ is always even
        \item To make an odd number, we must use an odd number of 15-cent coins
        \item Since $15c$ with $c$ odd gives odd contribution, we need $c \geq 1$ for odd totals
    \end{itemize}
    \item \textbf{Invariants}: The parity of the total is determined by the parity of $c$ (number of 15-cent coins)
    \item \textbf{Case Analysis}: Split into even numbers (use $c = 0, 2, 4, \ldots$) and odd numbers (use $c = 1, 3, 5, \ldots$)
    \item \textbf{Modular Arithmetic}: Consider values modulo 2, 3, 5, or 6 to determine representability
\end{itemize}

\subsection*{B. Domain-Specific Tactics}

\textbf{Number Theory:}
\begin{itemize}[leftmargin=*]
    \item \textbf{GCD/LCM}: $\gcd(6, 10, 15) = 1$, so eventually all integers are representable
    \item \textbf{Frobenius Number}: For two coprime integers $p$ and $q$, the Frobenius number is $pq - p - q = (p-1)(q-1) - 1$
    \item \textbf{Chicken McNugget Theorem}: Any integer $\geq (p-1)(q-1)$ can be written as $pa + qb$ with $a, b \geq 0$ when $\gcd(p,q) = 1$
    \item \textbf{Modular arithmetic}: Analyze which residue classes modulo 6 can be achieved
\end{itemize}

\subsection*{C. Crossover Tactics}
\begin{itemize}[leftmargin=*]
    \item \textbf{Number-Theoretic-Algebraic}: Diophantine equations with non-negative constraints
    \item \textbf{Combinatorial-Algebraic}: Counting which values can be represented
\end{itemize}
\end{tacticalbox}

\begin{techniquesbox}
\subsection*{A. High-Probability Techniques}
\begin{itemize}[leftmargin=*]
    \item \textbf{Primary Technique}: Frobenius/McNugget Theorem (M5 - Medium Probability)
    \item \textbf{Recognition Cues}: 
    \begin{itemize}
        \item Problem asks for ``largest value that cannot be made''
        \item Multiple coin denominations given
        \item Need non-negative integer combinations
        \item This is the classic signature of a Frobenius coin problem
    \end{itemize}
    \item \textbf{Application}: 
    \begin{enumerate}
        \item Recognize that $6a + 10b = 2(3a + 5b)$
        \item For $3a + 5b$: since $\gcd(3, 5) = 1$, by Chicken McNugget theorem, all integers $\geq (3-1)(5-1) = 8$ can be represented
        \item Therefore, all even integers $\geq 2 \cdot 8 = 16$ can be represented as $6a + 10b$
        \item For odd numbers, we need $c$ odd. With $c = 1$, we have $15 + 6a + 10b$, so all even $\geq 16$ become odd $\geq 31$
        \item This means all odd numbers $\geq 31$ are representable
        \item Therefore, all numbers $\geq 30$ are representable (30 is even, 31 is odd, and we have both)
        \item Check that 29 is not representable (29 is odd, so need $c$ odd; smallest is $c = 1$, giving $15 + 6a + 10b = 29$, so $6a + 10b = 14$. But 14 is not representable as $6a + 10b$)
    \end{enumerate}
\end{itemize}

\subsection*{B. Alternative Techniques}
\begin{itemize}[leftmargin=*]
    \item \textbf{Modular Arithmetic (H19)}: 
    \begin{itemize}
        \item Analyze which residue classes modulo 6 can be achieved
        \item $6 \equiv 0 \pmod{6}$, $10 \equiv 4 \pmod{6}$, $15 \equiv 3 \pmod{6}$
        \item All residues 0-5 mod 6 can eventually be achieved
        \item Find the largest value in each residue class that cannot be achieved
    \end{itemize}
    \item \textbf{Systematic Enumeration}: 
    \begin{itemize}
        \item List all representable values up to some threshold
        \item Identify gaps (non-representable values)
        \item Determine the largest gap
    \end{itemize}
    \item \textbf{Tabular Method (Visual)}:
    \begin{itemize}
        \item Arrange integers in rows of 6
        \item Mark representable values
        \item Identify the largest unmarked value
    \end{itemize}
    \item \textbf{Technique Integration}: Combine parity argument with Chicken McNugget theorem
\end{itemize}

\subsection*{C. Technique Selection Justification}
\begin{itemize}[leftmargin=*]
    \item \textbf{Why Frobenius/McNugget}: This is the standard technique for coin problems asking for largest non-representable value
    \item \textbf{Computational Simplicity}: By factoring out 2 and applying McNugget theorem, we reduce complexity
    \item \textbf{Reliability}: Well-established theorem with clear conditions
    \item \textbf{Time Efficiency}: Once recognized, solution takes 3-4 minutes
\end{itemize}
\end{techniquesbox}

\begin{verificationbox}
\subsection*{A. Verification Level Selection}
\begin{itemize}[leftmargin=*]
    \item \textbf{Chosen Level}: Level 5 - Thorough Verification (1-3 minutes)
    \item \textbf{Justification}: 
    \begin{itemize}
        \item Mid-to-upper band problem (P16)
        \item Requires verification that 29 is not representable
        \item Requires verification that 30 and above are representable
        \item Multiple methods available for cross-checking
        \item Confidence should be high (80-90\%) before moving on
    \end{itemize}
    \item \textbf{Time Budget}: 1.5-2 minutes for thorough verification
\end{itemize}

\subsection*{B. Verification Checklist}
\begin{itemize}[leftmargin=*]
    \item \textbf{Verify 29 is not representable}:
    \begin{itemize}
        \item 29 is odd, so need $c$ odd (at least $c = 1$)
        \item With $c = 1$: $6a + 10b = 29 - 15 = 14$
        \item Can we make 14 from $6a + 10b$? Test: $a = 0$: no (10 doesn't divide 14). $a = 1$: $10b = 8$ (no). $a = 2$: $10b = 2$ (no).
        \item 14 is not representable as $6a + 10b$ with $a, b \geq 0$
        \item Therefore, 29 is not representable
    \end{itemize}
    \item \textbf{Verify all numbers $\geq 30$ are representable}:
    \begin{itemize}
        \item Even numbers $\geq 30$: By McNugget theorem, all even $\geq 16$ can be made from $6a + 10b$
        \item Odd numbers $\geq 31$: $31 = 15 + 16$, $33 = 15 + 18$, $35 = 15 + 20$, etc. All representable
        \item Alternatively: all odd $\geq 31$ are $15 + (\text{even} \geq 16)$, which are all representable
    \end{itemize}
    \item \textbf{Compute digit sum}: $29 \to 2 + 9 = 11$
    \item \textbf{Check answer choices}: 11 corresponds to choice (D)
\end{itemize}

\subsection*{C. Alternative Verification Methods}
\begin{itemize}[leftmargin=*]
    \item \textbf{Method 1}: Explicit enumeration
    \begin{itemize}
        \item List representable values: 6, 10, 12, 15, 16, 18, 20, 21, 22, 24, 25, 26, 27, 28, 30, \ldots
        \item Non-representable: 1, 2, 3, 4, 5, 7, 8, 9, 11, 13, 14, 17, 19, 23, 29
        \item Largest: 29
    \end{itemize}
    \item \textbf{Method 2}: Modulo 6 table
    \begin{itemize}
        \item Residue 0 mod 6: 6, 12, 18, 24, 30, \ldots (all $\geq 6$ representable)
        \item Residue 1 mod 6: 31, \ldots (all $\geq 31$ representable using $15 + 16, 15 + 22, \ldots$)
        \item Residue 2 mod 6: 2, 8, 14, 20, 26, \ldots (20 and above representable)
        \item Residue 3 mod 6: 15, 21, 27, \ldots (all $\geq 15$ representable)
        \item Residue 4 mod 6: 4, 10, 16, 22, 28, \ldots (10 and above representable)
        \item Residue 5 mod 6: 5, 11, 17, 23, 29, 35, \ldots (35 and above representable, 29 is last gap)
    \end{itemize}
    \item \textbf{Method 3}: Direct verification of 30, 31, 32
    \begin{itemize}
        \item 30 = 6(5) + 10(0) + 15(0) = 30 \checkmark
        \item 31 = 6(1) + 10(1) + 15(1) = 6 + 10 + 15 = 31 \checkmark
        \item 32 = 6(2) + 10(2) + 15(0) = 12 + 20 = 32 \checkmark
        \item Pattern continues for all larger values
    \end{itemize}
\end{itemize}
\end{verificationbox}

\begin{solutionbox}
\subsection*{A. Step-by-Step Solution}
\begin{enumerate}[leftmargin=*]
    \item \textbf{Initial Setup}: 
    \begin{itemize}
        \item Let $a, b, c$ be non-negative integers representing the number of 6-cent, 10-cent, and 15-cent coins, respectively
        \item We seek the largest value $x$ such that $6a + 10b + 15c = x$ has no solution with $a, b, c \geq 0$
    \end{itemize}
    
    \item \textbf{Key Observation - Parity Analysis}: 
    \begin{itemize}
        \item Note that $6a + 10b = 2(3a + 5b)$ is always even
        \item Therefore, to make an odd number, we must have $c$ odd (use an odd number of 15-cent coins)
        \item To make an even number, we can use $c$ even (including $c = 0$)
    \end{itemize}
    
    \item \textbf{Apply Chicken McNugget Theorem to Even Numbers}:
    \begin{itemize}
        \item For even numbers, consider $6a + 10b = 2(3a + 5b)$
        \item Since $\gcd(3, 5) = 1$, by the Chicken McNugget theorem, all integers $\geq (3-1)(5-1) = 2 \cdot 4 = 8$ can be written as $3a + 5b$ with $a, b \geq 0$
        \item Therefore, all even integers $\geq 2 \cdot 8 = 16$ can be written as $6a + 10b$ with $a, b \geq 0$
    \end{itemize}
    
    \item \textbf{Determine Largest Non-Representable Even Number}:
    \begin{itemize}
        \item From step 3, all even numbers $\geq 16$ are representable
        \item Check even numbers $< 16$: 2, 4, 6, 8, 10, 12, 14
        \item Representable: 6 (1 coin of 6), 10 (1 coin of 10), 12 (2 coins of 6)
        \item Non-representable: 2, 4, 8, 14
        \item Largest non-representable even number: 14
    \end{itemize}
    
    \item \textbf{Apply Chicken McNugget Theorem to Odd Numbers}:
    \begin{itemize}
        \item For odd numbers, we need $c$ odd. The smallest is $c = 1$
        \item With $c = 1$: $6a + 10b + 15 = x$, so we need $6a + 10b = x - 15$
        \item For this to have a solution, we need $x - 15$ to be a representable even number
        \item From step 3, all even numbers $\geq 16$ are representable
        \item Therefore, all odd numbers $\geq 15 + 16 = 31$ are representable
    \end{itemize}
    
    \item \textbf{Determine Largest Non-Representable Odd Number}:
    \begin{itemize}
        \item From step 5, all odd numbers $\geq 31$ are representable
        \item Check odd numbers $< 31$: 1, 3, 5, 7, 9, 11, 13, 15, 17, 19, 21, 23, 25, 27, 29
        \item Representable with $c = 1$ (i.e., $15 + \text{even}$): 
        \begin{itemize}
            \item 15 = 15 (c=1, a=b=0)
            \item 21 = 15 + 6 (c=1, a=1, b=0)
            \item 25 = 15 + 10 (c=1, a=0, b=1)
            \item 27 = 15 + 12 (c=1, a=2, b=0)
        \end{itemize}
        \item For 29: $29 = 15 + 14$. But 14 is not representable as $6a + 10b$ (from step 4)
        \item Therefore, 29 is not representable
        \item Largest non-representable odd number: 29
    \end{itemize}
    
    \item \textbf{Conclusion - Largest Non-Representable Value}:
    \begin{itemize}
        \item Largest non-representable even: 14
        \item Largest non-representable odd: 29
        \item Overall largest non-representable: $x = 29$
    \end{itemize}
    
    \item \textbf{Verification that 29 is Indeed Non-Representable}:
    \begin{itemize}
        \item 29 is odd, so we need $c$ odd
        \item $c = 1$: $6a + 10b = 14$. Testing: $a = 0, 1, 2$: $10b = 14$ (no), $10b = 8$ (no), $10b = 2$ (no). Not representable.
        \item $c = 3$: $6a + 10b = 29 - 45 = -16 < 0$. Impossible.
        \item Higher odd $c$ give even more negative results.
        \item Therefore, 29 is not representable.
    \end{itemize}
    
    \item \textbf{Verification that All Values $\geq 30$ are Representable}:
    \begin{itemize}
        \item For even $n \geq 30$: Since all even $\geq 16$ are representable, this includes all even $\geq 30$
        \item For odd $n \geq 31$: Since all odd $\geq 31$ are representable (shown in step 5), this holds
        \item Spot check: $30 = 6 \cdot 5$, $31 = 6 + 10 + 15$, $32 = 6 \cdot 2 + 10 \cdot 2$. All work.
    \end{itemize}
    
    \item \textbf{Compute Digit Sum}:
    \begin{itemize}
        \item $x = 29$
        \item Sum of digits: $2 + 9 = 11$
    \end{itemize}
\end{enumerate}

\subsection*{B. Alternative Approaches}

\textbf{Method 2: Modular Arithmetic (Modulo 6 Analysis)}

Construct a modulo 6 table to track which residue classes are achievable:
\begin{itemize}
    \item $6 \equiv 0 \pmod{6}$
    \item $10 \equiv 4 \pmod{6}$
    \item $15 \equiv 3 \pmod{6}$
\end{itemize}

Using these, we can achieve:
\begin{itemize}
    \item Residue 0: Using $6a$ (all multiples of 6)
    \item Residue 1: Using $10 \cdot 2 + 15 \cdot 1 = 20 + 15 = 35 \equiv 1 \pmod{6}$
    \item Residue 2: Using $10 \cdot 1 = 10 \equiv 4 \pmod{6}$ and $10 \cdot 2 = 20 \equiv 2 \pmod{6}$
    \item Residue 3: Using $15 \cdot 1 = 15 \equiv 3 \pmod{6}$
    \item Residue 4: Using $10 \cdot 1 = 10 \equiv 4 \pmod{6}$
    \item Residue 5: Using $10 \cdot 2 + 15 \cdot 1 = 35 \equiv 5 \pmod{6}$. So residue 5 is achievable starting from 35.
\end{itemize}

The largest value in residue class 5 (mod 6) that cannot be achieved is 29.

\textbf{Method 3: Tabular Visual Method}

Arrange integers in rows of 6:
\begin{center}
\begin{tabular}{|c|c|c|c|c|c|}
\hline
1 & 2 & 3 & 4 & 5 & 6 \\
\hline
7 & 8 & 9 & 10 & 11 & 12 \\
\hline
13 & 14 & 15 & 16 & 17 & 18 \\
\hline
19 & 20 & 21 & 22 & 23 & 24 \\
\hline
25 & 26 & 27 & 28 & 29 & 30 \\
\hline
31 & 32 & 33 & 34 & 35 & 36 \\
\hline
\end{tabular}
\end{center}

Mark representable values (can verify by construction). The largest unmarked value is 29.

\textbf{Method 4: Apply Chicken McNugget Theorem Twice}

This is the most elegant approach:
\begin{align}
6a + 10b + 15c &= 2(3a + 5b) + 15c
\end{align}

For $3a + 5b$: Since $\gcd(3, 5) = 1$, the Frobenius number is $(3-1)(5-1) - 1 = 8 - 1 = 7$. So all integers $\geq 8$ can be written as $3a + 5b$.

Therefore:
\begin{align}
2(3a + 5b) + 15c &= 2t + 15c \quad \text{where } t \geq 8
\end{align}

This gives us $2t + 15c$ where $t \geq 8$. Now apply Chicken McNugget again:
\begin{itemize}
    \item We can make all values of the form $2t + 15c$ where $t \geq 8$ and $c \geq 0$
    \item This is equivalent to finding the Frobenius number for coefficients 2 and 15 with offset
    \item But wait, $\gcd(2, 15) = 1$, so eventually all integers are representable
    \item The Frobenius number for 2 and 15 is $(2-1)(15-1) - 1 = 14 - 1 = 13$
    \item But we have $2t$ where $t \geq 8$, so we start from $2 \cdot 8 = 16$
    \item All even $\geq 16$ are representable
    \item All odd $\geq 15 + 16 = 31$ are representable
    \item The largest non-representable is 29
\end{itemize}

\end{solutionbox}

\begin{competitionbox}
\subsection*{A. Time Management}
\begin{itemize}[leftmargin=*]
    \item \textbf{Estimated Time}: 4-5 minutes total
    \begin{itemize}
        \item Recognition and setup: 30 seconds
        \item Apply Chicken McNugget theorem: 2-3 minutes
        \item Verify 29 is non-representable: 1 minute
        \item Verify all $\geq 30$ are representable: 1 minute
        \item Compute digit sum and select answer: 15 seconds
    \end{itemize}
    \item \textbf{Time Checkpoints}: 
    \begin{itemize}
        \item At 1 minute: Should have recognized this as Frobenius/McNugget problem
        \item At 3 minutes: Should have applied theorem and identified 29 as candidate
        \item At 5 minutes: Should have verified and selected answer
    \end{itemize}
    \item \textbf{Priority Level}: Medium-High
    \begin{itemize}
        \item P16 is upper mid-band, worth attempting
        \item If Frobenius/McNugget theorem is familiar, this is very doable
        \item If not familiar, this could take 8-10 minutes
    \end{itemize}
\end{itemize}

\subsection*{B. Risk Assessment}
\begin{itemize}[leftmargin=*]
    \item \textbf{Confidence Level}: 85-90\% after thorough verification
    \begin{itemize}
        \item Multiple methods confirm 29 is correct
        \item Verification shows 29 is non-representable
        \item Verification shows all $\geq 30$ are representable
        \item Digit sum calculation is straightforward
    \end{itemize}
    \item \textbf{Abandonment Criteria}: 
    \begin{itemize}
        \item If after 3 minutes you haven't recognized the Frobenius structure, consider moving on
        \item Could return later with fresh perspective
        \item However, if you've identified 29 as a candidate, spend 1-2 more minutes verifying
    \end{itemize}
    \item \textbf{Flag for Review}: 
    \begin{itemize}
        \item No, if verification is thorough (multiple methods agree)
        \item Yes, if only one method used and still uncertain about 30 and above being representable
    \end{itemize}
\end{itemize}

\subsection*{C. Answer Format}
\begin{itemize}[leftmargin=*]
    \item Multiple choice: Select (D) 11
    \item Double-check: $2 + 9 = 11$ \checkmark
    \item Bubble carefully on answer sheet
\end{itemize}
\end{competitionbox}

\begin{keyinsightbox}
\textbf{Key Insight}: This is a Frobenius coin problem (Chicken McNugget theorem). The crucial observation is that $6a + 10b = 2(3a + 5b)$ produces only even numbers. By applying the Chicken McNugget theorem to the coprime pair (3, 5), we can determine that all even numbers $\geq 16$ are representable using just 6 and 10. For odd numbers, we need at least one 15-cent coin. The largest non-representable odd is $29 = 15 + 14$, where 14 itself is non-representable as $6a + 10b$. Thus, 29 is the Frobenius number for this system.
\end{keyinsightbox}

\begin{warningbox}
\textbf{Warning}: Don't forget to verify that 29 is actually non-representable! It's tempting to compute that all numbers $\geq 30$ are representable and assume 29 must be the answer, but you should explicitly check that $29 = 15 + 14$ fails because 14 cannot be made from $6a + 10b$. Also, be careful about the Chicken McNugget theorem conditions: it only applies to \textit{coprime} pairs. Here, $\gcd(6, 10) = 2 \neq 1$, so we factor out the 2 first to get the coprime pair (3, 5).
\end{warningbox}

\begin{tipbox}
\textbf{Pro Tip}: For Frobenius/coin problems with three denominations, try to reduce to a two-denomination problem using parity or modular arithmetic. In this case, recognizing that $6a + 10b$ produces only even numbers, and that 15 is needed for odd numbers, immediately splits the problem into two manageable cases: even (use 6 and 10) and odd (use 15 plus 6 and 10). This reduction is the key to efficient solution.

Additionally, memorize the Chicken McNugget theorem for two coprime integers $p$ and $q$: the largest non-representable value is $pq - p - q = (p-1)(q-1) - 1$, and all integers $\geq (p-1)(q-1)$ are representable. For $(3, 5)$, this gives Frobenius number 7 and threshold 8.
\end{tipbox}

\section*{Final Answer}

The largest value that cannot be purchased using 6-cent, 10-cent, and 15-cent coins with exact change is $x = 29$. The sum of the digits of 29 is $2 + 9 = 11$.

$$\boxed{\textbf{(D) } 11}$$

\section*{Confidence Assessment}
\begin{itemize}[leftmargin=*]
    \item \textbf{Confidence Level}: 90\% 
    \begin{itemize}
        \item Multiple verification methods all confirm 29
        \item Chicken McNugget theorem correctly applied
        \item Explicit verification that 29 is non-representable completed
        \item Explicit verification that 30, 31, 32, \ldots are representable completed
        \item Digit sum calculation is straightforward and verified
    \end{itemize}
    \item \textbf{Verification Applied}: Level 5 - Thorough Verification
    \begin{itemize}
        \item Primary method: Chicken McNugget theorem with parity analysis
        \item Alternative method: Modular arithmetic (mod 6 residue classes)
        \item Cross-check: Direct verification of representability for 29, 30, 31
    \end{itemize}
    \item \textbf{Flag for Review}: No
    \begin{itemize}
        \item High confidence after multiple verification methods
        \item Clear theoretical foundation (Chicken McNugget theorem)
        \item Computational verification confirms theoretical result
    \end{itemize}
    \item \textbf{Time Spent}: 
    \begin{itemize}
        \item Estimated: 4-5 minutes for solution + 1.5-2 minutes for verification = 5.5-7 minutes total
        \item This is appropriate for a P16 problem in the AMC 12
    \end{itemize}
\end{itemize}

\end{document}

